\Qhead{1.11}{Хаотические колебания в динамических системах}{Хаотические колебания в динамических системах. Понятие о хаотическом (странном) аттракторе. Возможные пути потери устойчивости регулярных колебаний и перехода к хаосу.}
\Subfile{1.11.1}{sections/01_oscillations/q11/s01}
\Subfile{1.11.2}{sections/01_oscillations/q11/s02}

\recap{Хаос --- ограниченное нерегулярное детерминированное движение с чувствительностью к начальному состоянию. Странный аттрактор связан с растяжением и складыванием. Переходы возможны через удвоения периода, перемежаемость и разрушение тора.}
\board{Три переменные диодного контура; аттрактор Лоренца; сечение Пуанкаре; бифуркационная диаграмма.}
\oral{Почему ни третья переменная, ни сплошной спектр сами по себе не доказывают хаос?}
\bait{Можно оставить окна периодичности: увеличение управляющего параметра не обязано монотонно усиливать хаос.}
\sources{Пятаков~\cite{pyatakov}, PDF-с.~69: диодный контур и дополнительная ёмкость; с.~70--71: Лоренц, Пуанкаре и признаки динамики. Rothman~\cite{rothman25,rothman28} --- сценарии перехода. Рисунки рассчитаны по указанным моделям, не являются экспериментом диодной схемы.}
